\documentclass[11pt]{article}
\usepackage[margin=1in]{geometry}
\usepackage{amsmath,amssymb,amsthm,mathtools}
\usepackage{microtype}
\usepackage[colorlinks=true,citecolor=blue,urlcolor=blue]{hyperref}
\newcommand{\Sp}{\mathcal S}
\newtheorem{theorem}{Theorem}
\title{A continuum of Schatten counterexamples with normal off-diagonal block}
\author{Partial answer to arXiv:2111.15180}
\date{13 August 2026}

\begin{document}
\maketitle

\begin{center}
\fbox{\parbox{0.88\linewidth}{\textbf{Status.} New proved partial result.
The proposed inequality fails for every $p>q_*=7.709989669514\ldots$.
The remaining finite-exponent range is not classified here.}}
\end{center}

\section{Question and result}

Bourin and Lee ask, after Proposition~3.4 of \cite{BL}, for which
$p\in[1,\infty]$ the inequality
\begin{equation}\label{eq:question}
 \left\|\begin{bmatrix}A&N\\N^*&B\end{bmatrix}\right\|_p
 \leq \|A+B\|_p
\end{equation}
holds for every positive block matrix when $N$ is normal.  Their proposition
proves \eqref{eq:question} for $p=2$; for $p=1$ both sides equal
$\operatorname{Tr}(A+B)$.

Define $q_*>1$ to be the second positive solution of
\begin{equation}\label{eq:qstar}
 \left(\frac{35}{32}\right)^q
 +2\left(\frac{29}{64}\right)^q=2.
\end{equation}
Numerically,
\[
 q_*=7.709989669514\ldots .
\]

\begin{theorem}\label{thm:main}
For every $p>q_*$, including $p=\infty$, there are
$A,B\in M_3^+$ and a unitary $N\in M_3$ such that
$\left[\begin{smallmatrix}A&N\\N^*&B\end{smallmatrix}\right]\geq0$ but
\eqref{eq:question} is false.
\end{theorem}

Thus the set of exponents sought in \cite{BL} is contained in
$[1,q_*]$.  The theorem strengthens a single operator-norm counterexample
to an explicit half-line of Schatten exponents.

\section{Exact construction}

Put
\[
 A=\operatorname{diag}\!\left(\frac25,\frac52,6\right),\qquad
 C=\operatorname{diag}\!\left(\frac52,\frac25,6\right),\qquad
 N=\begin{bmatrix}0&1&0\\0&0&1\\1&0&0\end{bmatrix},
\]
and set
\[
 B=N^*CN=\operatorname{diag}\!\left(6,\frac52,\frac25\right),
 \qquad H=\begin{bmatrix}A&N\\N^*&B\end{bmatrix}.
\]
Certainly $N$ is unitary, hence normal.  Congruence by
$\operatorname{diag}(I,N^*)$ gives
\[
 H\simeq \begin{bmatrix}A&I\\I&C\end{bmatrix}.
\]
Because $A$ and $C$ are diagonal, the latter matrix is permutation-similar
to the direct sum
\[
 \begin{bmatrix}2/5&1\\1&5/2\end{bmatrix}
 \oplus
 \begin{bmatrix}5/2&1\\1&2/5\end{bmatrix}
 \oplus
 \begin{bmatrix}6&1\\1&6\end{bmatrix}.
\]
The first two summands have determinant zero and the last is positive.
Consequently $H\geq0$, and its eigenvalues are exactly
\begin{equation}\label{eq:Hspec}
 7,\ 5,\ \frac{29}{10},\ \frac{29}{10},\ 0,\ 0.
\end{equation}
On the other hand,
\begin{equation}\label{eq:Sspec}
 A+B=\operatorname{diag}\!\left(\frac{32}{5},5,\frac{32}{5}\right).
\end{equation}

For finite $p\geq1$, equations \eqref{eq:Hspec}--\eqref{eq:Sspec} show
that $\|H\|_p>\|A+B\|_p$ precisely when
\[
 7^p+5^p+2(29/10)^p>2(32/5)^p+5^p,
\]
or, after cancellation and normalization, precisely when
\begin{equation}\label{eq:gineq}
 g(p):=\left(\frac{35}{32}\right)^p
 +2\left(\frac{29}{64}\right)^p>2.
\end{equation}

It remains only to justify the stated threshold.  We have $g(1)=2$ and
$g''(p)>0$.  Moreover $g'(1)<0$: indeed
\[
 35\log(35/32)+29\log(29/64)<35\frac3{32}-29\frac12<0,
\]
using $\log(1+x)<x$ and $64/29>2$ with $\log2>1/2$.
Finally $g(p)\to\infty$, since $35/32>1$.  Strict convexity therefore gives
exactly one further intersection $q_*>1$ with the level $2$, and
$g(p)>2$ exactly for $p>q_*$.  This proves the finite-$p$ assertion.  At
$p=\infty$, \eqref{eq:Hspec}--\eqref{eq:Sspec} give
$\|H\|_\infty=7>32/5=\|A+B\|_\infty$.

\section{Provenance and remaining gap}

Hayashi's Problem~3 and its following example \cite[PDF pp.~5--6]{Hay}
use the same cyclic-unitary reduction, with different diagonal entries, to
disprove the operator-norm inequality.  The rational choice above is tuned
so that one eigenvalue cancels and the entire Schatten comparison becomes
the two-exponential equation \eqref{eq:qstar}.

No claim is made for $1<p<2$ or $2<p\leq q_*$.  Endpoint interpolation does
not by itself compare the Schatten norms of two different matrices, and
tensoring or taking direct sums preserves the sign of a fixed $p$-norm
comparison.  These are substantive obstructions to upgrading this packet
to a full characterization.

\begin{thebibliography}{9}
\bibitem{BL}
J.-C. Bourin and E.-Y. Lee,
\emph{Eigenvalue inequalities for positive block matrices with the inradius
of the numerical range}, Int. J. Math. 33 (2022), 2250009;
arXiv:2111.15180, Proposition~3.4 and the question on PDF page~7.

\bibitem{Hay}
T. Hayashi, \emph{On a norm inequality for a positive block-matrix},
Linear Algebra Appl. 566 (2019), 86--97; arXiv:1808.00181.
\end{thebibliography}

\end{document}
